\section{Proof of Theorem~\ref{thm:intersect} (new version)}
\label{sec:proof-new}

To prove: $\mu_f^*$ is a geodesic current, extending $\mu_f$.

We have defined $\mu_f$ on the ring $\mathcal{R}$ generated by
right-handed boxes (defined to be half-open). We will construct an
associated geodesic current. (Loosely following
Hersonsky-Paulin~\cite{HP97:RigidityNegCurvedCone}?)

Work with oriented currents for simplicity of notation (no need for
$\Z/2$ quotient), although results ultimately apply to unoriented
currents.

As in Theorem~\ref{thm:caratheodory}, for $E \subset G^+(S)$ a subset
of the space of geodesics, we define
\[
    \mu_f^*(E) \coloneqq \inf \Bigl\{\, \sum_{i \in I} \mu_f(B_i) \Bigm| \text{$\{B_i\}$
          finite or countable cover of $E$ by right-handed boxes}\,\Bigr\}.
\]
Then $\mu_f^*$ is an outer measure.

\begin{lemma}
  The outer measure $\mu_f^*$ is a Borel measure.
\end{lemma}

\begin{proof}
  Since $\mathcal{R}$ generates Borel sets as a $\sigma$-algebra, it
  suffices to show that any right-handed box~$A \in \mathcal{R}$ is
  $\mu_f^*$-measurable.
  \DPTtodo[inline]{I started to write the proof, but it used no facts
    about $\mathcal{R}$ other than being a ring, which means it should
    be cited or proved in Section~\ref{sec:measure-thy}}
\end{proof}

It remains to prove that $\mu_f$ is countably additive, which is
equivalent to showing that $\mu_f^*$ and $\mu_f$ agree
on~$\mathcal{R}$. We introduce some notation.

For $B$ a right-handed box $[a_1,a_2) \times [b_1,b_2)$, let
$\overline{B} \coloneqq [a_1,a_2] \times [b_1,b_2]$ and
$B^\circ \coloneqq (a_1,a_2) \times (b_1,b_2)$ be the corresponding
closed and open boxes.

\begin{definition}
  For $B = [a_1,a_2) \times [b_1,b_2)$ a right-handed box and
  $\epsilon > 0$, let $V(B,\epsilon) = [a_1',a_2) \times [b_1',b_2)$
  be an element of $\mathcal{R}$ so that
  \begin{itemize}
  \item $a_1' < a_1$ and $b_1' < b_1$ (i.e., $V(B,\epsilon)$ extend
    $B$ slightly down and to the left)and
  \item $\mu_f(V(B,\epsilon) \setminus B) < \epsilon$.
  \end{itemize}
  It is possible to find such $a_1'$ and $b_1'$ by
  Proposition~\ref{prop:finalfamily1}.

  Also set $W(B,\epsilon)$ (similar neighborhood of upper/right edges
\end{definition}



%%% Local Variables:
%%% mode: latex
%%% TeX-master: "Intersections"
%%% End:
